\documentclass[10pt,a4paper]{amsart}
\usepackage{amsmath,amssymb,amsthm,graphicx}
% Retained minimal preamble for the Stacks Project chapter sites-cohomology.tex (Cohomology on Sites).
% Verbatim excerpt of the chapter preamble: Xy-pic setup (native preamble lines 14-18) and the native
% theorem declarations with the mathematical macros (native preamble lines 154-217). Copied mechanically;
% nothing added, removed or redefined. The wrapper class and the amsmath/amssymb/amsthm/graphicx base
% packages are supplied by the assembler, so the class/externaldocument part of the native preamble is omitted.
% For commutative diagrams we use Xy-pic
\usepackage[all]{xy}

% We use 2cell for 2-commutative diagrams.
\xyoption{2cell}
% Theorem environments.
%
\theoremstyle{plain}
\newtheorem{theorem}[subsection]{Theorem}
\newtheorem{proposition}[subsection]{Proposition}
\newtheorem{lemma}[subsection]{Lemma}

\theoremstyle{definition}
\newtheorem{definition}[subsection]{Definition}
\newtheorem{example}[subsection]{Example}
\newtheorem{exercise}[subsection]{Exercise}
\newtheorem{situation}[subsection]{Situation}

\theoremstyle{remark}
\newtheorem{remark}[subsection]{Remark}
\newtheorem{remarks}[subsection]{Remarks}

\numberwithin{equation}{subsection}

% Macros
%
\def\lim{\mathop{\mathrm{lim}}\nolimits}
\def\colim{\mathop{\mathrm{colim}}\nolimits}
\def\Spec{\mathop{\mathrm{Spec}}}
\def\Hom{\mathop{\mathrm{Hom}}\nolimits}
\def\Ext{\mathop{\mathrm{Ext}}\nolimits}
\def\SheafHom{\mathop{\mathcal{H}\!\mathit{om}}\nolimits}
\def\SheafExt{\mathop{\mathcal{E}\!\mathit{xt}}\nolimits}
\def\Sch{\mathit{Sch}}
\def\Mor{\mathop{\mathrm{Mor}}\nolimits}
\def\Ob{\mathop{\mathrm{Ob}}\nolimits}
\def\Sh{\mathop{\mathit{Sh}}\nolimits}
\def\NL{\mathop{N\!L}\nolimits}
\def\CH{\mathop{\mathrm{CH}}\nolimits}
\def\proetale{{pro\text{-}\acute{e}tale}}
\def\etale{{\acute{e}tale}}
\def\QCoh{\mathit{QCoh}}
\def\Ker{\mathop{\mathrm{Ker}}}
\def\Im{\mathop{\mathrm{Im}}}
\def\Coker{\mathop{\mathrm{Coker}}}
\def\Coim{\mathop{\mathrm{Coim}}}

% Boxtimes
%
\DeclareMathSymbol{\boxtimes}{\mathbin}{AMSa}{"02}

%
% Macros for moduli stacks/spaces
%
\def\QCohstack{\mathcal{QC}\!\mathit{oh}}
\def\Cohstack{\mathcal{C}\!\mathit{oh}}
\def\Spacesstack{\mathcal{S}\!\mathit{paces}}
\def\Quotfunctor{\mathrm{Quot}}
\def\Hilbfunctor{\mathrm{Hilb}}
\def\Curvesstack{\mathcal{C}\!\mathit{urves}}
\def\Polarizedstack{\mathcal{P}\!\mathit{olarized}}
\def\Complexesstack{\mathcal{C}\!\mathit{omplexes}}
% \Pic is the operator that assigns to X its picard group, usage \Pic(X)
% \Picardstack_{X/B} denotes the Picard stack of X over B
% \Picardfunctor_{X/B} denotes the Picard functor of X over B
\def\Pic{\mathop{\mathrm{Pic}}\nolimits}
\def\Picardstack{\mathcal{P}\!\mathit{ic}}
\def\Picardfunctor{\mathrm{Pic}}
\def\Deformationcategory{\mathcal{D}\!\mathit{ef}}

% --- appended 2026-10-12 (Chief): the project own definitions of the reference/slogan/history
% display environments, copied mechanically from the shared Stacks preamble lines 8-12
% (\usepackage{verbatim} plus the three \newenvironment lines, which the chapter preamble keeps
% verbatim from the project preamble). Nothing added, removed or redefined: these environments
% carry no mathematical content and the project itself typesets them as comments. Without these
% three lines a counted statement whose own text uses them fails with Environment slogan undefined.
\usepackage{verbatim}
\newenvironment{reference}{\comment}{\endcomment}
\newenvironment{slogan}{\comment}{\endcomment}
\newenvironment{history}{\comment}{\endcomment}
\usepackage{cleveref}
\crefname{originallemma}{original lemma}{original lemmas}
\Crefname{originallemma}{Original lemma}{Original lemmas}
\makeatletter
\newlabel{derived-lemma-adjoint-preserve-K-injectives}{{08BJ}{source}{}{sourcepointer-1}{}}
\newlabel{derived-lemma-adjoint-preserve-K-injectives@cref}{{[originallemma][0][]08BJ}{[1][0][]source}{}{}{}}
\makeatother
\crefname{originallemma}{original lemma}{original lemmas}
\Crefname{originallemma}{Original lemma}{Original lemmas}
\makeatletter
\newlabel{derived-lemma-leray-acyclicity}{{015E}{source}{}{sourcepointer-2}{}}
\newlabel{derived-lemma-leray-acyclicity@cref}{{[originallemma][0][]015E}{[1][0][]source}{}{}{}}
\makeatother
\crefname{originalreference}{original reference}{original references}
\Crefname{originalreference}{Original reference}{Original references}
\makeatletter
\newlabel{lemma-higher-direct-images}{{03Q8}{source}{}{sourcepointer-3}{}}
\newlabel{lemma-higher-direct-images@cref}{{[originalreference][0][]03Q8}{[1][0][]source}{}{}{}}
\makeatother
\crefname{originalreference}{original reference}{original references}
\Crefname{originalreference}{Original reference}{Original references}
\makeatletter
\newlabel{lemma-modules-abelian-unbounded}{{0D5Y}{source}{}{sourcepointer-4}{}}
\newlabel{lemma-modules-abelian-unbounded@cref}{{[originalreference][0][]0D5Y}{[1][0][]source}{}{}{}}
\makeatother
\crefname{originalreference}{original reference}{original references}
\Crefname{originalreference}{Original reference}{Original references}
\makeatletter
\newlabel{lemma-pullback-injective-pre-limp}{{0D6X}{source}{}{sourcepointer-5}{}}
\newlabel{lemma-pullback-injective-pre-limp@cref}{{[originalreference][0][]0D6X}{[1][0][]source}{}{}{}}
\makeatother
\crefname{originallemma}{original lemma}{original lemmas}
\Crefname{originallemma}{Original lemma}{Original lemmas}
\makeatletter
\newlabel{sites-modules-lemma-exactness-lower-shriek}{{04BH}{source}{}{sourcepointer-6}{}}
\newlabel{sites-modules-lemma-exactness-lower-shriek@cref}{{[originallemma][0][]04BH}{[1][0][]source}{}{}{}}
\makeatother
\crefname{originallemma}{original lemma}{original lemmas}
\Crefname{originallemma}{Original lemma}{Original lemmas}
\makeatletter
\newlabel{sites-modules-lemma-special-square-continuous}{{0FN4}{source}{}{sourcepointer-7}{}}
\newlabel{sites-modules-lemma-special-square-continuous@cref}{{[originallemma][0][]0FN4}{[1][0][]source}{}{}{}}
\makeatother
\crefname{originalsection}{original section}{original sections}
\Crefname{originalsection}{Original section}{Original sections}
\makeatletter
\newlabel{sites-section-cocontinuous-morphism-topoi}{{00XN}{source}{}{sourcepointer-8}{}}
\newlabel{sites-section-cocontinuous-morphism-topoi@cref}{{[originalsection][0][]00XN}{[1][0][]source}{}{}{}}
\makeatother
\crefname{originalsection}{original section}{original sections}
\Crefname{originalsection}{Original section}{Original sections}
\makeatletter
\newlabel{sites-section-morphism-sites}{{00X0}{source}{}{sourcepointer-9}{}}
\newlabel{sites-section-morphism-sites@cref}{{[originalsection][0][]00X0}{[1][0][]source}{}{}{}}
\makeatother
\title{Base change for the derived pushforward in a special continuous and cocontinuous square of ringed topoi}
\author{The Stacks Project Authors}
\date{stacks-project:sites-cohomology.tex:local-snapshot-20261011; GNU FDL 1.2 or later}
\begin{document}
\maketitle
\noindent\textit{Editorial note.} All delivered mathematical text is the author's own, copied from the pinned Stacks Project chapter file sites-cohomology.tex; no statement or proof step was rewritten, reordered or removed.\par\medskip
\noindent\textit{Editorial note.} The counted proof is the author's complete argument: the reduction to abelian sheaves via condition (5), the bounded below injective-resolution argument with the acyclicity check and Leray's acyclicity lemma, and the unbounded K-injective argument relying on the exactness of g'\_! in assumption (6).\par\medskip
\noindent\textit{Editorial note.} The label lemma-special-square-continuous names the counted statement, and the proof environment immediately following it in the source is delivered verbatim as its proof.\par\medskip
\noindent\textit{Editorial note.} The counted statement is carried by the author's own lemma environment in the chapter source.\par\medskip
\noindent\textit{Editorial note.} Pre-open gate: python operations/agent-authoring/tools/precard\_gate.py --source sites-cohomology.tex --label lemma-special-square-continuous --start 9594 --end 9681 --local-id 3555 returned RESULT=OK (L3555 free, occupancy FREE, public interval clear). Candidate taken from lane-s4/halfB-priority.json.\par\medskip
\noindent\textit{Editorial note.} The statement contains a footnote recording when condition (6) holds; the footnote is delivered verbatim with the statement.\par\medskip
\noindent\textit{Editorial note.} The statement and proof are copied from the pinned chapter file; the cited inputs are kept as external pointers with their canonical Stacks tags.\par\medskip
\noindent\textit{Allowed standard prerequisites (editorial).}\begin{itemize}
\item Pullback along a cocontinuous morphism with trivial pullback of the structure sheaf. Same chapter, the lemmas identifying g\textasciicircum{}* with g\textasciicircum{}\{-1\} when the pullback of the structure sheaf is the structure sheaf. Supplies the identification g\textasciicircum{}* = g\textasciicircum{}\{-1\} and (g')\textasciicircum{}* = (g')\textasciicircum{}\{-1\} from condition (5), so the derived functors reduce to underived pullbacks.
\item Vanishing of the higher cohomology of the pullback of an injective sheaf. Same chapter, the lemma that H\textasciicircum{}p(U', (g')\textasciicircum{}\{-1\}I\textasciicircum{}q) = 0 for p > 0 for an injective abelian sheaf I and a cocontinuous v'. Supplies the right acyclicity of the pullback of the terms of an injective resolution for the functor f'\_*.
\item Higher direct images as the sheaf associated to a presheaf of cohomology groups. Same chapter, the lemma identifying R\textasciicircum{}p f'\_* with the sheafification of V' maps to H\textasciicircum{}p(u'(V'), -). Supplies the translation of the vanishing of the cohomology groups over the objects U' into the acyclicity statement used by Leray's acyclicity lemma.
\item Leray's acyclicity lemma. Derived Categories, the lemma that a bounded below complex of objects acyclic for a left exact functor computes its right derived functor. Supplies the conclusion that f'\_*(g')\textasciicircum{}*I\textasciicircum{}bullet represents Rf'\_*(g')\textasciicircum{}\{-1\}K.
\item The commutativity of the underlying pushforward and pullback in the special square. Modules on Sites, the lemma computing the pushforward of the pullback along the square involved for sheaves of modules. Supplies the identity f'\_*(g')\textasciicircum{}\{-1\}I\textasciicircum{}bullet = g\textasciicircum{}\{-1\}f\_*I\textasciicircum{}bullet, which identifies the two sides of the base change map.
\item Adjoints of exact functors preserve K-injectivity. Derived Categories, the lemma that the right adjoint of an exact functor preserves K-injective complexes. Supplies, together with the exactness of g'\_! in condition (6), the K-injectivity of (g')\textasciicircum{}\{-1\}I\textasciicircum{}bullet in the unbounded case.
\end{itemize}\medskip
\ifdefined\hypertarget\hypertarget{sourcepointer-1}{}\fi
\noindent\textit{Source reference.} Original lemma 08BJ: sources/stacks-chow/tags.txt, source lines 10569--10569. Prerequisite pointer invoked at proof line(s) 9672; it names a standard fact of a sibling chapter that the counted argument uses, and every step of the counted proof built on it is retained verbatim in the delivered body.\par\medskip
\ifdefined\hypertarget\hypertarget{sourcepointer-2}{}\fi
\noindent\textit{Source reference.} Original lemma 015E: sources/stacks-chow/tags.txt, source lines 1505--1505. Prerequisite pointer invoked at proof line(s) 9659; it names a standard fact of a sibling chapter that the counted argument uses, and every step of the counted proof built on it is retained verbatim in the delivered body.\par\medskip
\ifdefined\hypertarget\hypertarget{sourcepointer-3}{}\fi
\noindent\textit{Source reference.} Original reference 03Q8: sources/stacks-sites-cohomology/pinned-original-source.tex, source lines 747--747. Prerequisite pointer invoked at proof line(s) 9656; it names a standard fact of the same chapter that the counted argument uses, and every step of the counted proof built on it is retained verbatim in the delivered body.\par\medskip
\ifdefined\hypertarget\hypertarget{sourcepointer-4}{}\fi
\noindent\textit{Source reference.} Original reference 0D5Y: sources/stacks-sites-cohomology/pinned-original-source.tex, source lines 4453--4453. Prerequisite pointer invoked at proof line(s) 9643; it names a standard fact of the same chapter that the counted argument uses, and every step of the counted proof built on it is retained verbatim in the delivered body.\par\medskip
\ifdefined\hypertarget\hypertarget{sourcepointer-5}{}\fi
\noindent\textit{Source reference.} Original reference 0D6X: sources/stacks-sites-cohomology/pinned-original-source.tex, source lines 9278--9278. Prerequisite pointer invoked at proof line(s) 9651; it names a standard fact of the same chapter that the counted argument uses, and every step of the counted proof built on it is retained verbatim in the delivered body.\par\medskip
\ifdefined\hypertarget\hypertarget{sourcepointer-6}{}\fi
\noindent\textit{Source reference.} Original lemma 04BH: sources/stacks-chow/tags.txt, source lines 5652--5652. Prerequisite pointer invoked at proof line(s) 9634; it names a standard fact of a sibling chapter that the counted argument uses, and every step of the counted proof built on it is retained verbatim in the delivered body.\par\medskip
\ifdefined\hypertarget\hypertarget{sourcepointer-7}{}\fi
\noindent\textit{Source reference.} Original lemma 0FN4: sources/stacks-chow/tags.txt, source lines 19554--19554. Prerequisite pointer invoked at proof line(s) 9663; it names a standard fact of a sibling chapter that the counted argument uses, and every step of the counted proof built on it is retained verbatim in the delivered body.\par\medskip
\ifdefined\hypertarget\hypertarget{sourcepointer-8}{}\fi
\noindent\textit{Source reference.} Original section 00XN: sources/stacks-chow/tags.txt, source lines 1225--1225. Prerequisite pointer invoked at proof line(s) 9617; it names a standard fact of a sibling chapter that the counted argument uses, and every step of the counted proof built on it is retained verbatim in the delivered body.\par\medskip
\ifdefined\hypertarget\hypertarget{sourcepointer-9}{}\fi
\noindent\textit{Source reference.} Original section 00X0: sources/stacks-chow/tags.txt, source lines 1202--1202. Prerequisite pointer invoked at proof line(s) 9616; it names a standard fact of a sibling chapter that the counted argument uses, and every step of the counted proof built on it is retained verbatim in the delivered body.\par\medskip
\begin{lemma}
\label{lemma-special-square-continuous}
Consider a commutative diagram
$$
\xymatrix{
(\Sh(\mathcal{C}'), \mathcal{O}_{\mathcal{C}'})
\ar[r]_{(g', (g')^\sharp)} \ar[d]_{(f', (f')^\sharp)} &
(\Sh(\mathcal{C}), \mathcal{O}_\mathcal{C}) \ar[d]^{(f, f^\sharp)} \\
(\Sh(\mathcal{D}'), \mathcal{O}_{\mathcal{D}'}) \ar[r]^{(g, g^\sharp)} &
(\Sh(\mathcal{D}), \mathcal{O}_\mathcal{D})
}
$$
of ringed topoi and suppose we have functors
$$
\xymatrix{
\mathcal{C}' \ar[r]_{v'} &
\mathcal{C} \\
\mathcal{D}' \ar[r]^v \ar[u]^{u'} &
\mathcal{D} \ar[u]_u
}
$$
such that (with notation as in
Sites, Sections \ref{sites-section-morphism-sites} and
\ref{sites-section-cocontinuous-morphism-topoi}) we have
\begin{enumerate}
\item $u$ and $u'$ are continuous and give rise to the morphisms
$f$ and $f'$,
\item $v$ and $v'$ are cocontinuous giving rise to the morphisms $g$ and $g'$,
\item $u \circ v = v' \circ u'$,
\item $v$ and $v'$ are continuous as well as cocontinuous, and
\item $g^{-1}\mathcal{O}_{\mathcal{D}} = \mathcal{O}_{\mathcal{D}'}$
and $(g')^{-1}\mathcal{O}_{\mathcal{C}} = \mathcal{O}_{\mathcal{C}'}$.
\end{enumerate}
Then $Rf'_* \circ (g')^* = g^* \circ Rf_*$ as functors
$D^+(\mathcal{O}_\mathcal{C}) \to D^+(\mathcal{O}_{\mathcal{D}'})$.
If in addition
\begin{enumerate}
\item[(6)] $g'_! : \textit{Ab}(\mathcal{C}') \to \textit{Ab}(\mathcal{C})$
is exact\footnote{Holds if fibre products and equalizers exist in
$\mathcal{C}'$ and $v'$ commutes with them, see
Modules on Sites, Lemma \ref{sites-modules-lemma-exactness-lower-shriek}.},
\end{enumerate}
then $Rf'_* \circ (g')^* = g^* \circ Rf_*$ as functors
$D(\mathcal{O}_\mathcal{C}) \to D(\mathcal{O}_{\mathcal{D}'})$.
\end{lemma}

\begin{proof}
We have $g^* = Lg^* = g^{-1}$ and $(g')^* = L(g')^* = (g')^{-1}$
by condition (5).
By Lemma \ref{lemma-modules-abelian-unbounded} it suffices
to prove the result on the derived category $D^+(\mathcal{C})$ or
$D(\mathcal{C})$ of abelian sheaves.

\medskip\noindent
Choose an object $K \in D^+(\mathcal{C})$.
Let $\mathcal{I}^\bullet$ be a bounded below complex of injective abelian
sheaves on $\mathcal{C}$ representing $K$. By
Lemma \ref{lemma-pullback-injective-pre-limp}
we see that $H^p(U', (g')^{-1}\mathcal{I}^q) = 0$ for
all $p > 0$ and any $q$ and any $U' \in \Ob(\mathcal{C}')$.
Recall that $R^pf'_*(g')^{-1}\mathcal{I}^q$ is the sheaf
associated to the presheaf $V' \mapsto H^p(u'(V'), (g')^{-1}\mathcal{I}^q)$,
see Lemma \ref{lemma-higher-direct-images}.
Thus we see that $(g')^{-1}\mathcal{I}^q$ is right acyclic
for the functor $f'_*$. By Leray's acyclicity lemma
(Derived Categories, Lemma \ref{derived-lemma-leray-acyclicity})
we find that $f'_*(g')^*\mathcal{I}^\bullet$
represents $Rf'_*(g')^{-1}K$.
By Modules on Sites, Lemma
\ref{sites-modules-lemma-special-square-continuous}
we find that $f'_*(g')^{-1}\mathcal{I}^\bullet = g^{-1}f_*\mathcal{I}^\bullet$.
Since $g^{-1}f_*\mathcal{I}^\bullet$ represents $g^{-1}Rf_*K$
we conclude.

\medskip\noindent
Choose an object $K \in D(\mathcal{C})$.
Let $\mathcal{I}^\bullet$ be a K-injective complex of abelian
sheaves on $\mathcal{C}$ representing $K$. By
Derived Categories, Lemma \ref{derived-lemma-adjoint-preserve-K-injectives}
and assumption (6) we find that $(g')^{-1}\mathcal{I}^\bullet$
is a K-injective complex of abelian sheaves on $\mathcal{C}'$.
By Modules on Sites, Lemma
\ref{sites-modules-lemma-special-square-continuous}
we find that $f'_*(g')^{-1}\mathcal{I}^\bullet = g^{-1}f_*\mathcal{I}^\bullet$.
Since $f_*\mathcal{I}^\bullet$ represents $Rf_*K$ and since
$f'_*(g')^{-1}\mathcal{I}^\bullet$ represents $Rf'_*(g')^{-1}K$
we conclude.
\end{proof}

\end{document}
